\scanpage{30}
\figure{\begin{tikzpicture}
\node[anchor=north west,inner sep=0pt,outer sep=0pt] (letter) at (0,0) {\begin{minipage}{0.98\linewidth}
\footnotesize
% letterhead
\begin{minipage}[t]{0.56\linewidth}
\centering\sffamily
UNIVERSITE DES SCIENCES ET TECHNIQUES\\
DU LANGUEDOC\\
\rule{0.55cm}{0.4pt}\\
INSTITUT DE MATHEMATIQUES\\
Place Eugène Bataillon\\
\unsure{34060|34050} MONTPELLIER CEDEX\\
\rule{0.65cm}{0.4pt}\\
\end{minipage}\hfill
\begin{minipage}[t]{0.40\linewidth}
Montpellier, le\\[2.4\baselineskip]
% Two detached compact marks to the left of the typed dateline; origin unassigned.
\makebox[0pt][r]{\begin{tikzpicture}[x=1pt,y=1pt,baseline=0pt]
\fill (0,0.4) ellipse[x radius=0.35pt,y radius=0.6pt,rotate=-12];
\fill (1.45,1.4) ellipse[x radius=0.4pt,y radius=0.65pt,rotate=-12];
\end{tikzpicture}\hspace{1.6em}}Villecun le 30.01.78\\
\end{minipage}

\smallskip
Monsieur Alexandre Grothendieck\\
Professeur\\

\hfill\begin{minipage}{0.49\linewidth}
à Monsieur le Ministre de\\
\unsure{\deleted{l'Education Nationale}}\\
\end{minipage}

\smallskip
Monsieurle Ministre,\\

\deleted{Je me permets de vous envoyer copie d'une lettre datée de ce même jour, adressée à M. Rouzaud, Administrateur Provisoire de l'U.S.T.L., dans laquelle je lui soumets certains des problèmes de fonctionnement posés aux Universités par l'article 21 de l'ordonnance du 2.11.1945.}

\smallskip
\deleted{Par le même courrier j'envoie également copie de la même lettre à M. Peyrefitte, Ministre de la Justice, en attirant son attention sur les difficultés créées par cette ordonnance. Celles-ci ne touchent certes pas, à mon sens, l'aspect le plus important des problèmes soulevés par l'existence même ce cette ordonnance. Par contre, par leur nature technique, elles me semblent bien du ressort des responsabilités' officielles d'un Administrateur ou d'un Ministre. C'est pourquoi il m'a semblé opportu\deleted{\unsure{m|n}}\unsure{n|m} de vous \makebox[0pt][l]{\raisebox{0.65ex}{\tiny ç}}soumettre, en votre qualité de Ministre de l'Education Nationale,\mbox{\begin{tikzpicture}[x=1ex,y=1ex,baseline=0pt]
\fill (0.18,0.06) ellipse [x radius=0.07,y radius=0.055];
\fill (0.42,0.38) ellipse [x radius=0.07,y radius=0.055];
\end{tikzpicture}} des problèmes de votre compétence. Il me semblerait extrêmement souhaitable que vous preniez position publiquement à leur sujet.}

\smallskip
\hspace*{0.18\linewidth}\deleted{Veuillez agréer, Monsieur le Ministre, l'expression de}\\
\deleted{ma considération très distinguée .}\\

\hfill\begin{minipage}{0.31\linewidth}
\handwritten{A Grothendieck}\\[0.7\baselineskip]
% The paired detached strokes lie below and just to the right of the initial A.
\hspace*{0.7em}\begin{tikzpicture}[x=1pt,y=1pt,baseline=0pt]
\fill (0,0) .. controls (0.3,0.3) and (0.25,0.8) .. (0.45,1.05)
 .. controls (0.95,1.55) and (1.45,1.25) .. (1.45,0.9)
 .. controls (1.2,0.55) and (0.65,0.5) .. (0.45,0.45)
 .. controls (0.35,0.1) and (0.2,-0.15) .. (0,0);
\fill (3.6,-0.05) .. controls (3.9,0.1) and (3.85,0.55) .. (3.95,0.75)
 .. controls (4.25,0.95) and (4.7,1.0) .. (4.9,1.1)
 .. controls (5.2,1.15) and (5.25,0.75) .. (4.95,0.6)
 .. controls (4.6,0.5) and (4.25,0.4) .. (4.15,0.2)
 .. controls (4.1,-0.1) and (3.8,-0.3) .. (3.6,-0.05);
\end{tikzpicture}
\end{minipage}\\
\end{minipage}};
\begin{scope}[x={(letter.north east)},y={(letter.south west)}]
% Unassigned angular dark mark at the upper right physical edge; no letter or numeral supplied.
\draw[black,line width=0.55pt,line cap=round] (1.035,-0.026)
 .. controls (1.034,-0.022) and (1.034,-0.020) .. (1.035,-0.018)
 .. controls (1.037,-0.017) and (1.040,-0.017) .. (1.041,-0.018)
 -- (1.041,-0.020);
\fill (1.040,-0.023) ellipse[x radius=0.4pt,y radius=0.25pt];
\draw[gray,very thin] (0.654,0.258) -- (0.644,0.235)
 .. controls (0.561,0.272) and (0.327,0.389) .. (0.401,0.404)
 .. controls (0.449,0.420) and (0.738,0.382) .. (0.651,0.456)
 .. controls (0.579,0.511) and (0.399,0.568) .. (0.422,0.624)
 .. controls (0.451,0.668) and (0.719,0.641) .. (0.654,0.713)
 .. controls (0.593,0.781) and (0.477,0.842) .. (0.508,0.994);
\draw[gray,very thin] (0.654,0.254)
 .. controls (0.588,0.297) and (0.461,0.379) .. (0.443,0.407)
 .. controls (0.414,0.445) and (0.537,0.422) .. (0.528,0.449)
 .. controls (0.517,0.477) and (0.402,0.533) .. (0.492,0.583)
 .. controls (0.590,0.638) and (0.508,0.674) .. (0.475,0.733)
 .. controls (0.428,0.806) and (0.414,0.817) .. (0.431,0.889);
\end{scope}
\end{tikzpicture}}
% archive sheet number: 29
